\chapter{Résultat}
%%%%%%%%%%%%%%%%%%%%%%%%%%%%
%%%%%%%%%%%%%%%%%%%%%%%%%%%%
L'évaluation financière des modèles repose sur une stratégie de \textit{value bet} à seuil adaptatif.
Pour chaque modèle $m$, un seuil optimal $\hat{s}_m$ est déterminé %sur le jeu de \textit{train} 
en maximisant
le profit net cumulé parmi les valeurs $s \in [0{,}40\,;\,0{,}85]$ par pas de $0{,}01$.
Un pari est déclenché sur le match $i$ si et seulement si deux conditions sont 
simultanément vérifiées :
\begin{enumerate}
  \item la probabilité prédite $\hat{p}_i$ dépasse le seuil du modèle : $\hat{p}_i \geq \hat{s}_m$ ;
  \item le pari constitue un \textit{value bet}, c'est-à-dire que la probabilité implicite du bookmaker
        est inférieure à celle du modèle : $\hat{p}_i > 1 / c_i$, où $c_i$ désigne la cote proposée.
\end{enumerate}
La mise dynamique suit une stratégie par paliers selon le niveau de confiance :
$10$\,€ lorsque $\hat{s}_m \leq \hat{p}_i < \hat{s}_m + 0{,}10$,
et $30$\,€ lorsque $\hat{p}_i \geq \hat{s}_m + 0{,}10$.
En l'absence de pari éligible, aucune mise n'est engagée.
Les seuils retenus après optimisation varient entre $0{,}62$ (MLP) et $0{,}81$ (Random Forest),
reflétant la dispersion des distributions de probabilités selon les architectures.
La performance financière est mesurée par le retour sur investissement (ROI) \ref{sec:roi}. 
Une autre stratégie est également mise en place : le principe est le même, 
seule la mise reste fixe à 10 €.
Nous pouvons ainsi comparer les deux stratégies de mise et ainsi chercher la meilleure.
%%%%%%%%%%%%%%%%%%%
%%%%%%%%%%%%%%%%%%%
\section{Performances financière}
\begin{figure}[H]
    \includegraphics[width=\textwidth]{../../plot/maj/ml_stat_4/bilan_financier.png}
    \caption{Comparaison des modèles}
    \label{fig:mlp_ml}
\end{figure}

\noindent Sur la figure \ref{fig:mlp_ml} nous avons comparé les modèles ML ainsi que les modèles Poisson. 
La métrique utilisée 
pour la comparaison est le ROI. Le Résultat "ROI Std 0.5" représente le seuil standard de 0.5, 
le résultat "ROI Seuil Opti" utilise donc le seuil optimal trouvé et la mise fixe. Enfin le 
"ROI Dyn Opti" utilise le seuil optimal ainsi que la stratégie de mise dynamique.

\noindent Nous pouvons donc observer que le modèle le plus rentable semble être le 
SVM à noyau RBF. De plus nous observons que la stratégie mêlant seuil optimal et 
mise dynamique permet d'obtenir un retour sur investissement de 5.71\% comparé 
à 4.9\% en utilisant simplement le seuil optimal. La valeur "VB" est présente pour les modèles 
ne présentant pas de seuil de confiance, le seul seuil retenu est celui du value bet.

\begin{table}[H]
\centering
\caption{Volume de paris engagés (nombre de matchs misés) par modèle et stratégie}
\label{tab:volume_paris}
\begin{tabular}{l|c|c|c}
\hline
\textbf{Modèle} & \textbf{Paris Standard} & \textbf{Paris Seuil Opti} & \textbf{Paris Dyn Opti} \\
\hline
\hline
\multicolumn{4}{l}{\textit{Modèles de Machine Learning}} \\
\hline
SVM RBF (base) & 1558 & 42 & 17 \\
RandomForest (base) & 1411 & 119 & 95 \\
MLP (Optimisé) & 1358 & 453 & 246 \\
SVM RBF (Optimisé) & 1445 & 627 & 311 \\
XGBoost (Base) & 1718 & 394 & 329 \\
LightGBM (Base) & 1738 & 478 & 387 \\
XGBoost (Optimisé) & 1769 & 697 & 429 \\
MLP (Base) & 1312 & 703 & 446 \\
LogisticRegression (Optimisé) & 1755 & 680 & 465 \\
RandomForest (Optimisé) & 1776 & 700 & 466 \\
LightGBM (Optimisé) & 1765 & 695 & 511 \\
LogisticRegression (Base) & 1784 & 723 & 522 \\
\hline
\hline
\multicolumn{4}{l}{\textit{Modèles Statistiques \& Références}} \\
\hline
Dixon-Coles (Optimisé + Temps) & 1399 & 1399 & 1399 \\
Dixon-Coles (Stats) & 1425 & 1425 & 1425 \\
Baseline (Toujours 1) & 3558 & 3558 & 3558 \\
\hline
\end{tabular}
\end{table}
Nous pouvons constater grâce à \ref{tab:volume_paris} que le SVM de base bien qu'aillant le 
meilleur ROI, présente le plus faible nombre de paris. Nous comprenons donc que ce modèle ne 
paris uniquement que lorsque la probabilité est relativement élevée. Le SVM optimisé semble 
montrer de bonnes performances financière tout en maintenant un nombre de paris supérieur 
à 300 et de bonnes performances prédictives détaillées ci-après.



\section{Performances prédictives}
\begin{figure}[H]
    \includegraphics[width=\textwidth]{../../plot/maj/ml_stat_4/metriq.png}
    \caption{Comparaison des métriques}
    \label{fig:mlp_ml_métrique}
\end{figure}


\noindent Sur ce tableau nous avons comparé les métriques détaillées dans \ref{chap:métriques}. Nous pouvons 
ainsi voir que tous les modèles dépassent la Baseline, cependant, le meilleur score f1 reste celui 
de la baseline. Cela s'explique par le fait que les modèles sont entraînés de sorte à optimiser le profit 
et non la qualité prédictive. 
Nous affichons également le score de "Recall" afin de mieux comprendre les résultats et les modèles.\\

Le principal objectif visant à dépasser la prédiction naïve est atteint, en 
effet sur les deux tableaux, nous observons bien que les modèles ML sont meilleurs. Nous pouvons 
également observer que le SVM (RBF) qui avait le meilleur ROI pocède un score-f1 très bas (0.038), mais 
avec une précision de 73\%. Il présente un recall très faible, de 1.95\%, il ignore donc plus de 98\% 
des bons paris pour ne pas prendre de risque. 
Nous observons que le meilleur modèles d'un point de vue prédictif est le SVM optimisé. Il 
présente le deuxième meilleur ROI et les meilleurs Performances prédictives. 



\section{Classement des modèles}
\begin{figure}[H]
        \centering
        \includegraphics[width=\textwidth]{../../plot/maj/ml_stat_4/cdd.png}
        \caption{Diagramme de différence critique (Log Loss)}
        \label{fig:cdd_ml_stat_logloss}
\end{figure}
\begin{figure}[H]
        \centering
        \includegraphics[width=\textwidth]{../../plot/maj/ml_stat_4/cdd_eqm.png}
        \caption{Diagramme de différence critique (Brier Score)}
        \label{fig:cdd_ml_stat_brier}
\end{figure}

\noindent Les figures \ref{fig:cdd_ml_stat_logloss} et \ref{fig:cdd_ml_stat_brier} 
permettent de classer les modèles et le bookmaker Bwin en utilisant la métrique de 
log loss, respectivement le Brier score.\\

Les "crossbar" regroupent les modèles ne présentant pas de différence 
significative, ainsi plusieurs groupes peuvent être 
construits. Le plus important reste le groupe des équivalents du bookmaker, 
[Bwin, SVM (opti), RandomForest (opti), LogistiqueRegression (opti), LightGBM (opti)].



\section{Conclusion et Perspectives}

En conclusion, ce stage aura permis de mettre en place une méthodologie de la 
collecte des données jusqu'à la prédiction et l'évaluation des modèles. 

Les résultats obtenus indiquent clairement que les modèles de Machine Learning 
surpassent l'approche statistique classique, tant en termes de capacité prédictive 
que de rentabilité. Ces performances mettent toutefois en évidence un élément crucial : 
la rentabilité à long terme dépend directement de la sélectivité du modèle. 
Ce dernier doit restreindre son volume de paris via un seuil optimal pour ne pas 
prendre de risque inutile et perdre de l'argent.

Plusieurs pistes restent néanmoins ouvertes pour approfondir ce 
travail. Tout d'abord, l'extension de la base de données constituerait une amélioration 
naturelle : intégrer davantage de championnats (compétitions européennes secondaires, 
voire d'autres confédérations) et un historique plus long permettrait d'augmenter la 
taille de l'échantillon d'entraînement et d'évaluer la généralisation des modèles à des 
contextes footballistiques différents. De même, l'intégration de données en direct, telles 
que les statistiques de possession, les changements tactiques, les blessures des joueurs ou 
la composition des équipes, pourrait significativement enrichir l'information disponible 
et ouvrir la voie à une prédiction dynamique, actualisée tout au long de la rencontre.

Toutefois, il reste relativement compliqué au vu de la visualisation section \ref{sec:visu} 
de prédire avec certitude les matchs de foot. En effet, grace aux méthodes de visualisation 
nous voyons clairement que les données sont relativement peu séparables, ce qui explique 
nos résultats plutôt faible en terme de performances prédictive et financière. 
La méthodologie de sélection de variables pourrait elle aussi être affinée, par 
exemple en explorant des approches d'embedding ou des méthodes de réduction de dimension 
non linéaires en amont de la sélection, afin de capturer des interactions plus complexes 
entre les descripteurs. 

